\documentclass[12pt,letterpaper]{report}

\usepackage[utf8]{inputenc}
\usepackage[spanish]{babel}
\usepackage{amsmath}
\usepackage{amsfonts}
\usepackage{amssymb}
\usepackage{float}
\usepackage{makeidx}
\usepackage{graphicx}
\usepackage{color}
%\usepackage{kpfonts} % change font
%\usepackage{mathptmx} % change font
\usepackage{multirow}
%\usepackage[numbers]{natbib}
\usepackage[authoryear,comma]{natbib}
\usepackage{amsmath,amssymb,amsfonts,latexsym,cancel}
\usepackage[left=3cm,right=3cm,top=2cm,bottom=2cm]{geometry}
\usepackage{subfigure}
\usepackage{url}
\usepackage[breaklinks,colorlinks=true,linkcolor=black,citecolor=blue,urlcolor=black]{hyperref}
\usepackage{graphicx}
\usepackage{makecell}
\usepackage{enumitem, hyperref}
\usepackage[table]{xcolor}
\usepackage{longtable}
\usepackage{booktabs}
\usepackage{appendix}
\usepackage{rotating}
\usepackage[spanish,onelanguage,ruled,vlined,lined,resetcount]{algorithm2e}
\SetKw{Break}{finalizar ciclo;}
\usepackage[paper=portrait, pagesize]{typearea}
\usepackage{fancyhdr} % para poner el título en cada capitulo
%\usepackage{rotating} % para rotar las paginas
\usepackage{pdflscape}

% genera el mes de forma automatica
\usepackage{datetime}
\makeatletter
\newdateformat{mifecha}{La Habana, \monthname[\THEMONTH] \THEYEAR}
\renewcommand{\monthnamespanish}[1][\month]{
	\@orgargctr=#1\relax
	\ifcase\@orgargctr
	\PackageError{datetime}{Invalid Month number \the\@orgargctr}{
		Month number should go from 1 to 12
		}
		\or Enero
		\or Febrero
		\or Marzo
		\or Abril
		\or Mayo
		\or Junio
		\or Julio
		\or Agosto
		\or Septiembre
		\or Octubre
		\or Noviembre
		\or Diciembre
		\else \PackageError{datetime}{Invalid Month number \the\@orgargctr}{
			Month number should go from 1 to 12
		}
		\fi
		}


% permite crear listas dentro de lasa tablas
\newcolumntype{e}[1]{%--- Enumerated cells ---
	>{\minipage[t]{\linewidth}%
		\NoHyper%                Hyperref adds a vertical space
		\let\\\tabularnewline
		\enumerate
		\addtolength{\rightskip}{0pt plus 50pt}% for raggedright
		\setlength{\itemsep}{-\parsep}}%
	p{#1}%
	<{\@finalstrut\@arstrutbox\endenumerate
		\endNoHyper
		\endminipage}}

\newcolumntype{i}[1]{%--- Itemized cells ---
	>{\minipage[t]{\linewidth}%
		\let\\\tabularnewline
		\itemize
		\addtolength{\rightskip}{0pt plus 50pt}%
		\setlength{\itemsep}{-\parsep}}%
	p{#1}%
	<{\@finalstrut\@arstrutbox\enditemize\endminipage}}
\makeatother

% esto es para que funcionen los anexos incluyendo la ñ
\makeatletter \renewcommand*{\numberline}[1]{% 
	\hb@xt@\@tempdima{% 
		#1% 
		\protected@edef\@temp@num{#1}% 
		\ifx\@temp@num\@empty\else .\fi \hfil }% 
	} \makeatother


\newenvironment{dedication}
{\clearpage           % we want a new page
	\thispagestyle{empty}% no header and footer
	\vspace*{\stretch{.5}}% some space at the top 
	\itshape             % the text is in italics
	\raggedleft          % flush to the right margin
}
{\par % end the paragraph
	\vspace{\stretch{3}} % space at bottom is three times that at the top
	\clearpage           % finish off the page
}


\title{Título tesis}

\makeindex 

\renewcommand{\baselinestretch}{1.2}
\begin{document}
	\pagenumbering{roman}	
	\renewcommand{\listtablename}{Índice de tablas}
	\renewcommand{\tablename}{Tabla}
	\renewcommand{\bibname}{Referencias bibliográficas}
	
%	esto es para que las listas anidadas salgan con numeros
	\renewcommand{\labelenumii}{\arabic{enumi}.\arabic{enumii}}
	\renewcommand{\labelenumiii}{\arabic{enumi}.\arabic{enumii}.\arabic{enumiii}}
	\renewcommand{\labelenumiv}{\arabic{enumi}.\arabic{enumii}.\arabic{enumiii}.\arabic{enumiv}}
	
	
	\pagestyle{empty}	
	\markboth{}{} 	
	\thispagestyle{empty} 
	
 	\begin{figure}
	\centering
	\includegraphics[width=0.5\textwidth]{figuras/membrete-cujae-centrado.png}
\end{figure}

	\begin{center}
	
	\large{Trabajo de diploma para optar por el título de Ingeniería en Informática}

	\vspace{2cm}
	
	\LARGE{\textbf{Aquí va el Título tesis}}
	
	\vspace{2cm}
	
	\large{
	Autor:\\
	\emph{Aquí va el Nombre completo del estudiante}
	\vspace{0.5cm}
	
	Tutores:\\
	\emph{Aquí va el Nombre completo del primer tutor}\\
	Dr. C., Ing. Joaquín Danilo Pina Amargós 
	}
	
	\vspace{4cm}
	
	\small{\mifecha\today}
	
\end{center}	



	{\Large \textbf{Declaración de autoría} \vspace{.3cm}}

Declaro que soy el único autor de este trabajo y autorizo a la Facultad de Ingeniería Informática de la Universidad Tecnológica de la Habana para que haga el uso que estime pertinente de este trabajo. Para que así conste, firmo la presente el 7 de Marzo del 2025.

\vspace{3cm}

\begingroup	

\setlength{\tabcolsep}{10pt} % Default value: 6pt
\renewcommand{\arraystretch}{0.5} % Default value: 1
%\hspace{-1.2cm}
\centering
\begin{tabular}{c c}
	
%		\includegraphics[width=0.2\textwidth]{figuras/firma-autor.png}
%	 & \includegraphics[width=0.4\textwidth]{figuras/firma-tutor.png} \\
	Carlos Daniel Vilaseca Illanit  & Dra. C. Raisa Socorro Llanes \\
	\noindent\rule{6cm}{0.4pt} & \noindent\rule{6cm}{0.4pt} \\
	(Nombre y Apellidos del Diplomante) &  	(Nombres y apellidos del Tutor 1)	
\end{tabular}
%firma_jpinas
\vspace{3cm}

\endgroup


	\newgeometry{left=3cm,right=3cm,top=2cm,bottom=.5cm} 

{\Large \textbf{Opinión del tutor} \vspace{.3cm}}


Tutores:

\renewcommand{\arraystretch}{0.3} % Default value: 1
\begin{tabular}{l c}
%nombre tutor & \includegraphics[width=0.12\textwidth]{figuras/firmatutor.png} \\
\end{tabular}\\
%\vspace{.2cm}
\begin{center}
Cujae, La Habana, Cuba, Diciembre 2022
\end{center}

\renewcommand{\arraystretch}{1} % Default value: 1
\restoregeometry

	\begin{dedication}
esto es opcional
\end{dedication}


	{\Large \textbf{Agradecimientos} \vspace{.3cm}}

Opcional



	\section*{Resumen}
Aquí el resumen

\begin{description}
	\item[Palabras claves:]{la lista de palabras.}
\end{description}
%\end{abstract}



	\begin{abstract}
The challenge of forming teams of appropriate software projects is of great importance for software development companies, given that the team plays an important role in the success or failure of a project. The present work consists of developing an application to manage the data corresponding to instances of said problem, the application is organized in three fundamental parts for the management of your data, the first where all the sets to be taken into account to form teams are taken into account, the second where the relationships established between said sets are found, and in the third part is where the user is allowed to select the data that they wish to have in the instance to be managed. It is important to note that the data corresponding to the relationships can be generated randomly following statistical distributions, besides the intancias can be generated from randomly choosing the sets that will belong to it. On the other hand it is possible to generate simple solutions to the problem in a random way and see its evaluation. The application also has a database where all the managed information is stored, which allows access to previously inserted instances to clone, modify or simply view them. This database can be exported or imported in Excel format. Python was used as programming language, Pycharm as development IDE, PyQt as visual development framework, JasperSoft for reporting and PostgreSQL as database manager.
\end{abstract}
	

	\markboth{}{}
	\pagestyle{plain}	
	
	\tableofcontents	
	\pagebreak	
	
	\listoftables
	\pagebreak
	
	\listoffigures	
 	\clearpage 
	
	\pagebreak	
	\pagenumbering{arabic}
	
	\phantomsection
	\addcontentsline{toc}{chapter}{Introducción}
	\pagestyle{fancy} 
		
	\fancyhf{}
	\lhead[\thepage]{\textbf{Introducción}}
	\rhead[\textbf{Introducción}]{\thepage}
	
	\chapter*{Introducción}
Describir el objeto de estudio con las definiciones fundamentales. No olvidar las referencias bibliográficas, aquí tienen un ejemplo \citep{pina2019incorporation}.\\

Las figuras deben referenciarse en el texto, así como lo muestra este ejemplo, ver Figura \ref{fig:figCUJAE}.

\begin{figure}[H] %la opción H indica al compilador LaTeX que posicione la figura lo más cerca posible de este lugar.
\centering
  \includegraphics[width=0.5\linewidth]{figuras/membrete-cujae-centrado.png}
  \caption{El título de la figura debe estar acorde con su contenido.}
  \label{fig:figCUJAE} %incluir el label permite referenciarla en cualquier parte del documento.
\end{figure}

Se deben utilizar siempre los mismos términos para referirse a los mismos conceptos y no olvidar de definir los términos que son claves en el campo de acción, o sea, la propuesta de la tesis.\\

Describir la Situación problemática. Problema. Objetivo general. Objetivos específicos. Tareas. Beneficios. Breve resumen del contenido de la tesis.	
	
	\fancyhf{}
	\lhead[\thepage]{\leftmark}
	\rhead[\rightmark]{\thepage}
	\renewcommand{\chaptermark}[1]{\markboth{\chaptername \, \thechapter. #1}{}}
	
	\chapter{Marco teórico referencial de la investigación}\label{chap:1}

Incluir un resumen del contenido del capítulo en un párrafo.\\

\section{Título del epígrafe}

Texto... Es una buena práctica culminar (o iniciar) cada epígrafe con una referencia al que sigue (o precede) para dar fluidez a su lectura y no se vea como un \emph{copia y pega} sin ninguna relación.

\subsection{Título del subepígrafe}

Texto...

\subsection{Título del subepígrafe}

Texto...

\section{Título del epígrafe}

Texto...

\subsection{Título del subepígrafe}

Texto...

\subsection{Título del subepígrafe}

Texto...

\section{Conclusiones parciales}

% Cada conclusión tiene que estar sustentada en el cuerpo del capítulo.

Una vez terminado el capítulo se arriban a las siguientes conclusiones:

\begin{enumerate}
	\setlength\itemsep{0em}
	\item Una conclusión necesaria aquí son los requisitos principales que debe cumplir la solución propuesta.
	\item Otra conclusión es la inexistencia de una solución que brinde cumplimiento a los requisitos planteados.
	\item Finalmente, cuáles son las tecnologías seleccionadas y su justificación
\end{enumerate}

\pagebreak

	\chapter{Solución propuesta}\label{chap:2}

Incluir un resumen del contenido del capítulo en un párrafo.\\

\section{Título del epígrafe}

Texto... Es una buena práctica culminar (o iniciar) cada epígrafe con una referencia al que sigue (o precede) para dar fluidez a su lectura y no se vea como un \emph{copia y pega} sin ninguna relación.

\subsection{Título del subepígrafe}

Texto...

\subsection{Título del subepígrafe}

Texto...

\section{Título del epígrafe}

Texto...

\subsection{Título del subepígrafe}

Texto...

\subsection{Título del subepígrafe}

Texto...

\section{Conclusiones parciales}

% Cada conclusión tiene que estar sustentada en el cuerpo del capítulo.

Una vez terminado el capítulo se arriban a las siguientes conclusiones:

\begin{enumerate}
	\setlength\itemsep{0em}
	\item Unas conclusiones necesarias aquí son los aspectos tecnológicos que se tuvieron en cuenta para cumplir con los requisitos principales.
\end{enumerate}
\pagebreak


	\chapter{Validación de la solución}\label{chap:3}
Incluir un resumen del contenido del capítulo en un párrafo.\\

\section{Título del epígrafe}

Texto... Es una buena práctica culminar (o iniciar) cada epígrafe con una referencia al que sigue (o precede) para dar fluidez a su lectura y no se vea como un \emph{copia y pega} sin ninguna relación.

\subsection{Título del subepígrafe}

Texto...

\subsection{Título del subepígrafe}

Texto...

\section{Título del epígrafe}

Texto...

\subsection{Título del subepígrafe}

Texto...

\subsection{Título del subepígrafe}

Texto...

\section{Conclusiones parciales}

% Cada conclusión tiene que estar sustentada en el cuerpo del capítulo.

Una vez terminado el capítulo se arriban a las siguientes conclusiones:

\begin{enumerate}
	\setlength\itemsep{0em}
	\item Una conclusión se debe centrar en lo que se permite verificar con el diseño de las pruebas.
	\item Y por supuesto, el resultado como tal de las pruebas, limitaciones y beneficios de la solución propuesta respecto a las existentes.	
\end{enumerate}
\pagebreak

	
	
    
    \cleardoublepage
    \phantomsection
	\addcontentsline{toc}{chapter}{Conclusiones}
    \fancyhf{}
    \lhead[\thepage]{\textbf{Conclusiones}}
    \rhead[\textbf{Conclusiones}]{\thepage}
    \chapter*{Conclusiones}
Al finalizar este trabajo se arriban a las siguientes conclusiones:
\begin{enumerate}
	\item Estas conclusiones son un resumen e integración de las conclusiones parciales de todos los capítulos, pero nunca son un copia y pega de estas.
\end{enumerate}
    
    \cleardoublepage
    \phantomsection
	\addcontentsline{toc}{chapter}{Recomendaciones}
	\fancyhf{}
	\lhead[\thepage]{\textbf{Recomendaciones}}
	\rhead[\textbf{Recomendaciones}]{\thepage}
	\chapter*{Recomendaciones}
Con vistas a darle continuidad a este trabajo se recomienda: 
\begin{enumerate}
	\item Aquí se deben reflejar los aspectos que no dio tiempo incorporar en la solución propuesta u otros que brindan al lector la posibilidad de continuar el trabajo.
\end{enumerate}
	
	\cleardoublepage
	\phantomsection
	\addcontentsline{toc}{chapter}{Referencias bibliográficas}
	\fancyhf{}
	\lhead[\thepage]{\textbf{Referencias bibliográficas}}
	\rhead[\textbf{Referencias bibliográficas}]{\thepage}
	\bibliographystyle{authordate1}
	\bibliography{bib_tesis}

	
	\cleardoublepage
	\phantomsection
	\addcontentsline{toc}{chapter}{Anexos}
	\fancyhf{}
	\lhead[\thepage]{\textbf{Anexos}}
	\rhead[\textbf{Anexos}]{\thepage}
	\appendix
\clearpage{\renewcommand{\appendixname}{Anexo}

\chapter{Aquí va el Título del anexo}
Texto y figuras del anexo...

\chapter{Aquí va el Título del otro anexo}
Texto y figuras del anexo...

\end{document}